   \begin{center}
        {\large \textbf{Вариант №22 \\
        }}
    \end{center}

    \textbf{\underline{Цель работы:} изучить основные понятия теории формальных языков и грамматик.
    } 

    \begin{enumerate}
        \item \hyperref[sec:1]{Найти терминальную цепочку $\alpha$, $|\alpha| > 10$, для которой существует не менее двух левых выводов в заданной КС-грамматике (см. варианты
        заданий). Записать различные левые выводы этой цепочки. Построить деревья вывода. Определить последовательности правил, применяемые при этих
        выводах.}
        \item \hyperref[sec:2]{Написать программу, которая определяет, можно ли применить заданную последовательность правил при левом выводе терминальной
        цепочки в заданной КС-грамматике, формирует левый вывод и линейную скобочную форму дерева вывода. Обработать программой последовательности правил, полученные в
        п.1.
        Примечание. Если к нетерминалу А в процессе вывода применяется
        правило с номером n, то в выводе и в линейной скобочной форме дерева вывода после нетерминала А должен быть символ с кодом n.}
        \item \hyperref[sec:3]{Найти последовательность правил р, |p| > 10, которую можно применить при произвольном выводе терминальной цепочки, но нельзя
        применить при левом или правом выводе в заданной КС-грамматике
        (см. варианты заданий).
        Записать вывод v, в процессе которого применяется последовательность правил р. Построить дерево вывода.
        Записать левый и правый выводы, эквивалентные выводу v.}
        \item \hyperref[sec:4]{Написать программу, которая определяет, можно ли применить заданную последовательность правил р при выводе терминальной цепочки в заданной КС-грамматике и формирует линейную скобочную
        форму дерева вывода. Если последовательность правил р можно применить при выводе v терминальной цепочки, то программа должна вывести последовательность правил, применяемую при левом выводе,
        эквивалентном выводу v. Обработать программой последовательность правил, найденную в п.3.
        Примечание. Если к нетерминалу А в процессе вывода применяется
        правило с номером n, то в линейной скобочной форме дерева вывода
        после нетерминала А должен быть символ с кодом n.}
    \end{enumerate}